\documentclass{article}
\usepackage[margin=1in]{geometry}
\usepackage{amsmath,amssymb,amsthm}
\usepackage{hyperref}
\usepackage{tcolorbox}
\tcbuselibrary{breakable}

\newtcolorbox{accept}[1][]{colback=green!5,colframe=green!60!black,
  fonttitle=\bfseries,title={Accepted: #1},breakable}
\newtcolorbox{partial}[1][]{colback=yellow!5,colframe=orange!60!black,
  fonttitle=\bfseries,title={Partially Accepted: #1},breakable}
\newtcolorbox{correct}[1][]{colback=blue!5,colframe=blue!60!black,
  fonttitle=\bfseries,title={Author Correction: #1},breakable}

\title{Response to Reviewer op46 (claude-4.6-opus)}
\date{April 2026}
\begin{document}
\maketitle

\section{Overall Assessment}

op46 provides the most detailed and technically sophisticated review, scoring
5/10 and identifying several genuine issues of varying severity. We accept or
partially accept every major finding, as each is mathematically valid. The most
impactful contributions are: (1) the corrected Lemma~1 proof avoiding the
$J_\varphi(z_i) = I_d$ error; (2) the identification of the torus geodesic
bug; (3) the topological obstruction theorem.

\section{Weaknesses -- Response}

\subsection{W1: Lemma 1 Proof -- Normal Coordinate Jacobian}

\begin{accept}[Accepted -- proof will be revised]
op46 correctly identifies that the claim ``$J_\varphi(z_i) = I_d$ at the base
point'' is imprecise. In geodesic normal coordinates, the \emph{metric} at the
origin is $\delta_{kl}$ (NC2), but the Jacobian of the coordinate chart
$\varphi$ from Euclidean coordinates to normal coordinates is
$J_\varphi(z_i) = G(z_i)^{-1/2}Q$ for some orthogonal $Q$, which equals $I_d$
only when $G(z_i) = I_d$.

The corrected proof provided by op46 (via geodesic initial velocity $v =
\exp_{z_i}^{-1}(z_j)$ and Taylor expansion of the ODE) is rigorous and gives
the same final bound. Specifically:
\[
  z_j = z_i + v - \tfrac{1}{2}\Gamma^k_{lm}(z_i)v^l v^m \hat{e}_k + O(\|v\|^3)
  \implies v = \Delta z_{ij} + O(\|\Delta z\|^2)
\]
and the cross term $\langle \Delta z, c \rangle_{G(z_i)} = O(K_{\max}\|\Delta z\|^3)$
(via Christoffel identity), giving the same final bound $|d_R - w_{ij}| \leq
K_{\max}r^3/6$.

\textbf{Action}: Rewrite Step 3 of Lemma 1 proof. Replace the claim
``$J_\varphi(z_i) = I_d$'' with the geodesic velocity argument. The final
bound is unchanged.
\end{accept}

\subsection{W2: Assumption A4' -- Circular?}

\begin{accept}[Accepted -- will be stated as capacity condition]
op46 correctly identifies that the derivable bound is $O(\delta_{\mathrm{rec}}/r_n
+ r_n)$ and A4' postulates a strictly stronger bound without deriving it.

The resolution is the capacity condition $\delta_{\mathrm{rec}} = O(r_n^2)$:
under this condition, $\delta_{\mathrm{rec}}/r_n = O(r_n)$, so A4' follows
automatically from the finite-difference argument. The paper's Remark is already
correct; we will promote this to a formal Proposition.

\textbf{Action}: Add a Proposition after the Remark stating:
``If $\delta_{\mathrm{rec}} = O(r_n^2)$, then A4' holds with
$\|G^* - g\|_{\mathrm{op}} = O(r_n)$, and Theorem~2 becomes:
$\mathbb{E}|\dR^* - d_M| \leq C_1 r_n + \frac{K_{\max}}{6}r_n^3 + C_3 e^{-\lambda T/2}$.''
\end{accept}

\subsection{W3: E-step Not Descent -- Convergence Fragile}

\begin{partial}[Accepted as limitation -- not resolvable in general]
op46 correctly notes that the convergence analysis is purely local. We
acknowledge this explicitly in the paper. The empirical evidence (graphs
converge within 15--98 iterations from Euclidean initialization) provides
practical justification.

The fixed-point existence proof from reviewer Snn46 (Brouwer/finite-set
argument) will be added to address the existence gap. We cannot prove global
convergence for neural networks in full generality -- this is labeled as
Open Problem~1.
\end{partial}

\subsection{W4: Curvature Proxy Not a Consistent Estimator}

\begin{correct}[Critical theoretical correction]
op46 raises that $\hat{\kappa}$ converges to $C(M,u,v) \cdot K_M$ where $C$
is geometry-dependent. This is correct.

\textbf{The resolution (from Snn46, the most important theoretical finding):}
$\hat{\kappa} \to 2/R$ holds if and only if the decoder is isometrically
calibrated (pullback metric = true metric). For an uncalibrated decoder with
scale factor $\alpha$, $\hat{\kappa} \to (2/R\alpha)\sqrt{3/G}$.

This means $\hat{\kappa}$ is simultaneously:
\begin{enumerate}
  \item A curvature detector: $\hat{\kappa} = 0 \Leftrightarrow K = 0$ locally.
  \item An isometry quality metric: $\hat{\kappa} \to 2\sqrt{K}$ \emph{at and
    only at} the isometric fixed point.
\end{enumerate}

The experimental target $\hat{\kappa}^* = 2.0$ (for $S^2(R=1)$) is the
isometrically-calibrated value. The baseline's $\hat{\kappa} = 2.849$ overshoots
because Euclidean KNN edges introduce spurious curvature. SCR-VAE's $1.442$
is moving toward the isometric target.

\textbf{Action}: Add Proposition (from Snn46 bridge) showing $\hat{\kappa} \to
2/R$ iff decoder is isometrically calibrated. This transforms the experimental
comparison from ``curvature accuracy'' to ``simultaneous curvature detection
AND isometry quality measurement.''
\end{correct}

\subsection{W5: Experimental Gains Modest}

\begin{partial}[Accepted but reframed]
The gains (4.4\% Euclidean MAE, 34\% curvature) are modest in absolute terms.
However, two factors reframe this:
\begin{enumerate}
  \item The torus experiment uses the WRONG manifold (standard round torus
    instead of Clifford flat torus). After fixing, the curvature detection
    on the torus should improve dramatically.
  \item The sphere gains are limited by $n = 3000$ data points -- the
    theoretical rate $O(r_n) = O((\log n / n)^{1/2})$ limits how much the
    self-consistent iteration can improve for finite $n$.
  \item The training budget comparison needs revision: both models now use
    100k epochs (fair), but the SCR-VAE uses iterative E-steps which add
    graph construction overhead. The total wall-clock time is 13s/iter $\times$
    100 iter = 22 min vs. 16s for baseline -- comparable per-epoch cost.
\end{enumerate}
\end{partial}

\subsection{W6: No Comparison to Existing Methods}

\begin{partial}[Acknowledged -- deferred to revision]
op46 correctly notes no comparison to existing Riemannian VAE methods. This is
a valid weakness. Arvanitidis et al.\ (2018) treat the metric as a learnable
function of the latent code and compute geodesics numerically -- a fundamentally
different approach. Adding this comparison is important but requires significant
implementation effort. We defer this to a revision.
\end{partial}

\section{Critical Code Bugs -- Response}

\subsection{CRITICAL BUG: Torus Geodesic Distances}

\begin{accept}[Accepted -- highest priority fix]
op46 is completely correct. The function \texttt{compute\_true\_geodesic\_distances}
uses the flat torus metric $d = \sqrt{(R\Delta\theta)^2 + (r\Delta\phi)^2}$,
but the data is generated on the \textbf{standard embedded torus in $\mathbb{R}^3$}
with metric $ds^2 = (R + r\cos\phi)^2 d\theta^2 + r^2 d\phi^2$.

For $R=2, r=1$: at the inner equator ($\phi=\pi$), the true metric factor is
$(R-r)^2 = 1$ while the flat formula uses $R^2 = 4$ -- a factor of 4 error in
the $\theta$-direction metric.

\textbf{Fix}: Switch to the Clifford torus in $\mathbb{R}^4$ for which:
\begin{itemize}
  \item The induced metric IS flat ($K = 0$ everywhere).
  \item The geodesic formula IS $\sqrt{(R\Delta\theta)^2 + (r\Delta\phi)^2}$.
  \item The embedding matrix becomes $A \in \mathbb{R}^{G \times 4}$.
\end{itemize}
This is both a code fix AND makes the theory-experiment alignment exact.
\end{accept}

\subsection{CRITICAL: Torus Curvature Ground Truth Wrong}

\begin{accept}[Accepted]
The paper claims $\hat{\kappa}^* = 0$ (flat torus) but the data manifold has
nonzero $K(\phi)$. After switching to the Clifford torus, both the data AND
the curvature prediction will be correct: $K = 0$ everywhere, so
$\hat{\kappa}^* = 0$ genuinely.

The observed $\hat{\kappa} = 2.212$ in the current torus experiment is partially
measuring the genuine nonzero curvature of the standard embedded torus, not
a failure of the method.
\end{accept}

\subsection{Table: Code--Theory Alignment}

\begin{accept}[Accepted: symmetrized distances and distance clipping need documentation]
Both the symmetrized Riemannian distances $(w_{fwd} + w_{bwd})/2$ and the
distance clipping at $5 \times \mathrm{median}(w_{ij})$ are implemented but
not in the paper. Both will be added to the E-step description with theoretical
justification (both are within Lemma 1's $O(K_{\max} r^3)$ error budget).
\end{accept}

\section{Topological Obstruction Theorem}

\begin{accept}[Accepted -- will be added to paper]
The topological obstruction theorem (Theorem 1 in op46's Section 4.2) is a
genuine mathematical contribution that formally explains the torus performance
gap. We will add a version of this theorem to the paper (Section on
``Sphere vs.\ Torus''), stating:

\textit{Any $C^1$ decoder $f: \mathbb{R}^2 \to \mathbb{R}^G$ approximating
a compact manifold $M$ with $\pi_1(M) \neq 0$ must satisfy at least one of:
(a) incomplete coverage, (b) metric degeneracy (folds), or (c) curvature
concentration with total absolute curvature $\geq 4\pi$.}

This turns the torus ``failure'' into a theoretical \textbf{result}: the method
correctly identifies the fundamental impossibility of isometric decoding from
a simply-connected latent space.
\end{accept}

\section{Bridging the Gaps: Evaluation of op46's Proofs}

\subsection{Corrected Lemma 1 Proof (Section 3.1 of op46)}

\textbf{Mathematical evaluation}: op46's proof is correct. The key insight is
using the geodesic ODE and Taylor expansion:
\[
  z_j = z_i + v - \tfrac{1}{2}\Gamma^k_{lm}(z_i)v^l v^m \hat{e}_k + O(\|v\|^3)
\]
to get $v = \Delta z + O(\|\Delta z\|^2)$, then showing the cross-term
$\langle \Delta z, c \rangle_{G}$ involves $G_{kp}\Gamma^k_{lm}(\Delta z)^l(\Delta z)^m$
which is $O(K_{\max}\|\Delta z\|^3)$ by the Christoffel symbol--metric identity.
The final bound $|d_R - w_{ij}| = O(K_{\max} r^3)$ with coefficient $K_{\max}/6$
follows from the Bertrand-Puiseux expansion.

\textbf{Decision}: We will use this proof structure for the revised Lemma 1.

\subsection{Topological Obstruction (Section 4.2 of op46)}

\textbf{Mathematical evaluation}: The proof sketch is valid. The key argument
is: for a surjection $f: \mathbb{R}^2 \to T^2$, non-injectivity is forced by
simple connectivity, giving metric degeneracy (case b). If injectivity is
enforced by restricting to a domain, the image cannot be all of $T^2$ (case a).
The curvature concentration (case c) via Gauss-Bonnet is correct in spirit but
the bound $4\pi$ needs more careful justification for smooth folds.

\textbf{Decision}: We will add a simplified version (cases a/b only, which are
cleaner) to the paper.

\subsection{Quantifying Flat vs. Curved Torus Discrepancy (Section 4.3 of op46)}

\textbf{Mathematical evaluation}: The $50\%$ relative error estimate is correct.
For $R=2, r=1$: at $\phi=\pi$, the true $\theta$-metric coefficient is
$(R-r)^2 = 1$ while the flat formula uses $R^2 = 4$, giving a factor-of-2
discrepancy. This completely invalidates quantitative torus isometry comparisons.

\section{Summary of Accepted vs. Disputed Findings}

\begin{tabular}{lll}
\hline
Finding & Decision & Priority \\
\hline
W1: Lemma 1 proof error & Accept & Medium \\
W2: A4' not derived & Accept (capacity condition) & High \\
W3: Local convergence only & Accept as limitation & Low \\
W4: Curvature not consistent estimator & Accept + reframe as isometry metric & High \\
W5: Modest gains & Partial (reframe) & Medium \\
W6: No Riemannian VAE baselines & Deferred & Low \\
Critical Bug: torus geodesic & Accept (Clifford fix) & Critical \\
Critical Bug: torus curvature & Accept & Critical \\
Lemma 1 corrected proof & Accept & Medium \\
Topological obstruction theorem & Accept & High \\
\hline
\end{tabular}

\end{document}
